\section{Introduction}

\subsection{Problem Context}

Mobile applications ask users to make privacy decisions while the evidence needed to evaluate those decisions remains fragmented. A permission declaration states capability, not necessarily use; a runtime observation can indicate activity, but not purpose, lawful basis, necessity, or proportionality. This distinction is especially important for health-related applications, where personal data can fall within the GDPR's special-category regime. The Regulation requires legal and organisational analysis in addition to technical safeguards, including the principles in Article 5 and data protection by design and by default in Article 25 \cite{eurlex2016,edpb2019}.

Privacy Lens addresses a narrower engineering problem: how can an ordinary Android application present available permission-audit evidence without transforming incomplete telemetry into a false legal conclusion? The application collects no remote evidence, has no account or advertising layer, and stores findings locally. It exposes whether a finding came from device-visible observations or a synthetic demonstration, displays missing-evidence limitations, and prepares questions for human review.

The work began as a GDPR-oriented prototype. Revision 9 retains GDPR as the evaluated legal objective but removes GDPR from the core decision mechanism. Legal mappings are loaded through replaceable regulation packs. This separation makes regional extension possible while preventing the interface, repository, and Android bridge from silently inheriting one jurisdiction's terminology.

\subsection{Research Questions}

The thesis answers four questions:

\begin{enumerate}
    \item How can Android permission-audit signals be transformed into deterministic, explainable review findings while unsafe or incomplete inputs fail closed?
    \item How can evidence provenance, temporal scope, and missing legal context be communicated so that a technical warning is not mistaken for a legal verdict?
    \item How can jurisdiction-specific rules be isolated behind a stable regulation-pack contract without weakening traceability?
    \item What claims are supported by deterministic tests, release builds, and bounded physical-device evaluation, and which claims remain unverified?
\end{enumerate}

\subsection{Objectives and Contributions}

The first objective is a reproducible decision pipeline. Every finding retains its permission category, observed count, time window, source, regulation-pack identifier, rule reference, risk level, rationale, and missing-context fields. The second objective is calibrated communication. The interface distinguishes \emph{needs review}, \emph{evidence gap}, and \emph{no technical concern}; none is labelled compliant or non-compliant. The third objective is architectural replaceability. A registry loads immutable rule-pack definitions and the compliance engine consumes their shared contract. The fourth objective is deployable engineering: a stable Android package name, target SDK 36, release APK and Android App Bundle, explicit privacy policy, owner-action checklist, branded assets, and physical-device evidence.

The contribution is therefore not an automated GDPR judge. It is an accountability-support prototype that makes a bounded technical inference reconstructable and leaves legal conclusions to an appropriately authorised human reviewer.

\subsection{Scope and Claim Boundary}

The evaluated data path covers permission summary records available to the application or generated by its controlled simulator. Stock Android restricts ordinary applications from reading unrestricted activity belonging to other applications. Absence of an observation is consequently an evidence gap, not proof that an access did not occur. The current rule thresholds are research baselines; they are not statutory limits. The included EU GDPR pack is a traceable mapping aid, not legal advice.

Evaluation covers source compilation, deterministic logical tests, adversarial boundary tests, release assembly, manifest inspection, installation, launch, interaction, and short-run observation on one OPPO PERM00 device. It does not establish multi-OEM compatibility, 24-hour scheduler reliability, battery performance, population accuracy, accessibility conformance, independent legal validity, or application-store approval.

\subsection{Thesis Structure}

Chapter 2 establishes the legal, Android, HCI, and software-architecture foundations. Chapter 3 derives requirements and the evidence-bounded architecture. Chapter 4 documents the implemented application and regulation-pack mechanism. Chapter 5 reports reproducible evaluation evidence. Chapter 6 analyses validity, usability, deployment, and legal limitations. Chapter 7 answers the research questions and defines the next validation programme.
